\documentclass[../main.tex]{subfiles}
\graphicspath{{\subfix{../img/}}{\subfix{../architectures images/}}}
\begin{document}
\chapter{Sp\'{e}cification des Exigences}

\section{Introduction}
Avant d'\'{e}crire une ligne de code, il faut savoir ce qu'on construit et pour qui. Ce chapitre r\'{e}pond \`{a} ces questions : acteurs du syst\`{e}me, besoins fonctionnels et non fonctionnels, cas d'utilisation, product backlog et mod\`{e}le de donn\'{e}es.

\section{Sp\'{e}cification des exigences}

\subsection{Identification des Acteurs}
TrustGym distingue quatre acteurs, avec des p\'{e}rim\`{e}tres d'acc\`{e}s strictement d\'{e}finis :

\begin{itemize}
  \item \textbf{SuperAdmin} : responsable de la gestion globale de la plateforme \acrshort{saas}. Il g\`{e}re les partenaires, les salles de sport, les modules \acrshort{saas} (activation/d\'{e}sactivation via feature flags), le registre des portes d'acc\`{e}s, les bundles de modules et les utilisateurs de la plateforme.

  \item \textbf{Admin} : administrateur d'une salle de sport sp\'{e}cifique. Il dispose d'un acc\`{e}s complet \`{a} la gestion de sa salle : membres, coachs, abonnements, contr\^{o}le d'acc\`{e}s, boutique, finance, centre \acrshort{kpi} et assistant \acrshort{ai}.

  \item \textbf{Caissier} : op\'{e}rateur au niveau de la caisse. Il prend en charge l'inscription des membres, l'encaissement des paiements, l'enr\^{o}lement aux abonnements et les ventes en boutique. Il dispose d'un acc\`{e}s restreint par rapport \`{a} l'Admin. En raison du taux de rotation \'{e}lev\'{e} des caissiers dans le march\'{e} tunisien des salles de sport (3 \`{a} 6 mois en moyenne), des restrictions strictes ont \'{e}t\'{e} mises en place :
  \begin{itemize}
    \item \textbf{Interdiction de suppression} : le Caissier ne peut supprimer aucune donn\'{e}e (membres, factures, abonnements, paiements).
    \item \textbf{Interdiction de modification des param\`{e}tres} : il ne peut pas acc\'{e}der aux r\'{e}glages de la salle, aux formules d'abonnement, aux conventions ni aux salles.
    \item \textbf{Interdiction d'acc\`{e}s aux rapports financiers} : les rapports de ventes, les journaux de connexion et les portefeuilles pr\'{e}pay\'{e}s sont r\'{e}serv\'{e}s \`{a} l'Admin.
  \end{itemize}

  \item \textbf{Membre} (application mobile) : adh\'{e}rent de la salle de sport. Il acc\`{e}de \`{a} son tableau de bord, consulte les s\'{e}ances, r\'{e}serve des cr\'{e}neaux, visualise l'occupation en temps r\'{e}el et g\`{e}re son profil via l'application mobile d\'{e}di\'{e}e.
\end{itemize}

\subsection{Besoins Fonctionnels}
Le p\'{e}rim\`{e}tre fonctionnel par acteur, synth\'{e}tis\'{e} (le diagramme de cas d'utilisation global de la figure~\ref{fig:use_case_global} en donne la vue compl\`{e}te) :

\begin{itemize}
  \item \textbf{SuperAdmin} : partenaires, salles, modules \acrshort{saas} (\acrshort{ai}, Boutique, Coachs, \acrshort{kpi}, Finance, Acc\`{e}s), registre des portes, bundles de modules, utilisateurs plateforme.
  \item \textbf{Admin} : tableau de bord \acrshort{kpi}, \acrshort{crud} membres, abonnements, contr\^{o}le d'acc\`{e}s (entr\'{e}es/sorties, cr\'{e}neaux, occupation temps r\'{e}el), coachs et s\'{e}ances, boutique (produits/fournisseurs/stock/transferts), finance (wallet, facturation \acrshort{b2b}, caisse), assistant \acrshort{ai}/\acrshort{rag}.
  \item \textbf{Caissier} : inscription des membres, enr\^{o}lement aux abonnements, encaissement (esp\`{e}ces/ch\`{e}que), acc\`{e}s manuels, ventes boutique, rechargement wallet --- sans suppression ni acc\`{e}s aux rapports.
  \item \textbf{Membre} (mobile) : inscription multi-\'{e}tapes (e-mail, \acrshort{otp}, identit\'{e}, mot de passe), tableau de bord, calendrier des s\'{e}ances, r\'{e}servation des cr\'{e}neaux RESERVATION, statistiques (pr\'{e}sence, occupation, heures de pointe), profil.
\end{itemize}

\subsection{Besoins non fonctionnels}
Au-del\`{a} des fonctionnalit\'{e}s, la qualit\'{e} d'un syst\`{e}me se mesure \`{a} ses contraintes non fonctionnelles \cite{functional_req} :

\begin{itemize}
  \item \textbf{Performance} : utilisation du \acrshort{ssr} via Next.js pour un rendu rapide c\^{o}t\'{e} serveur. Mise en cache intelligente gr\^{a}ce \`{a} TanStack React Query avec un \texttt{staleTime} de 2 minutes et un cache de 15 minutes, r\'{e}duisant significativement les requ\^{e}tes r\'{e}seau.

  \item \textbf{Scalabilit\'{e}} : architecture modulaire NestJS permettant l'ajout de nouveaux modules sans impact sur les existants. D\'{e}ploiement sur \acrshort{vps} OVH avec PM2 pour une orchestration fiable des processus en production.

  \item \textbf{S\'{e}curit\'{e}} : double syst\`{e}me d'authentification \acrshort{jwt} (Admin/Membre). Protection des en-t\^{e}tes HTTP via Helmet. Configuration \texttt{CORS} restrictive. Contr\^{o}le d'acc\`{e}s bas\'{e} sur les r\^{o}les (\acrshort{rbac}). \texttt{ModuleGuard} pour la gestion des fonctionnalit\'{e}s \acrshort{saas} par salle.

  \item \textbf{Convivialit\'{e}} : internationalisation compl\`{e}te avec support du fran\c{c}ais, de l'arabe et de l'anglais, incluant la gestion automatique du mode RTL pour l'arabe. Interface responsive gr\^{a}ce \`{a} Tailwind \acrshort{css}.

  \item \textbf{Temps r\'{e}el} : communication en temps r\'{e}el via WebSocket \cite{rfc6455} (Socket.IO) pour la mise \`{a} jour instantan\'{e}e du taux d'occupation. M\'{e}canisme de polling toutes les 30 secondes en cas de d\'{e}connexion du WebSocket comme solution de repli.

  \item \textbf{Multi-tenant} : architecture multi-locataire \cite{bezemer2010multitenancy} avec isolation des donn\'{e}es par partenaire et par salle, permettant \`{a} plusieurs cha\^{i}nes de salles de sport de coexister sur la m\^{e}me instance de la plateforme.
\end{itemize}

\section{Analyse des besoins}

\subsection{Diagramme de cas d'utilisation global}
La figure~\ref{fig:use_case_global} synth\'{e}tise, d'un seul coup d'\oe{}il, ce que chaque acteur peut faire dans TrustGym.

\begin{figure}[H]
  \centering
  \resizebox{!}{0.8\textheight}{%
  \begin{tikzpicture}[
    link/.style={semithick}
  ]

  % System boundary : single column of use cases (crossing-free)
  \node[systembox, minimum width=5cm, minimum height=26cm] (system) at (5,-11.875) {};
  \node[uctitle] at (system.north) {Syst\`{e}me TrustGym};

  % --- Use cases : single column at x=5, grouped by actor (top -> bottom) ---
  % SuperAdmin
  \node[usecase] (uc1)  at (5,   0.00) {G\'{e}rer les partenaires};
  \node[usecase] (uc2)  at (5,  -1.25) {G\'{e}rer les salles};
  \node[usecase] (uc3)  at (5,  -2.50) {G\'{e}rer les modules SaaS};
  \node[usecase] (uc4)  at (5,  -3.75) {G\'{e}rer les portes};
  % Admin
  \node[usecase] (uc6)  at (5,  -5.00) {G\'{e}rer les abonnements};
  \node[usecase] (uc8)  at (5,  -6.25) {G\'{e}rer les coachs};
  \node[usecase] (uc9)  at (5,  -7.50) {Planifier les s\'{e}ances};
  \node[usecase] (uc10) at (5,  -8.75) {G\'{e}rer la boutique};
  \node[usecase] (uc11) at (5, -10.00) {G\'{e}rer la finance};
  \node[usecase] (uc12) at (5, -11.25) {Consulter les KPI};
  \node[usecase] (uc13) at (5, -12.50) {Utiliser l'assistant IA};
  % Admin + Caissier (shared)
  \node[usecase] (uc5)  at (5, -13.75) {G\'{e}rer les membres};
  \node[usecase] (uc7)  at (5, -15.00) {Contr\^{o}ler les acc\`{e}s};
  % Caissier
  \node[usecase] (uc14) at (5, -16.25) {Inscrire un membre};
  \node[usecase] (uc15) at (5, -17.50) {Encaisser un paiement};
  \node[usecase] (uc16) at (5, -18.75) {Vendre un produit};
  % Membre
  \node[usecase] (uc17) at (5, -20.00) {S'inscrire (mobile)};
  \node[usecase] (uc18) at (5, -21.25) {Consulter les s\'{e}ances};
  \node[usecase] (uc19) at (5, -22.50) {R\'{e}server une s\'{e}ance};
  \node[usecase] (uc20) at (5, -23.75) {Consulter l'occupation};

  % --- Actors (left), each centred on its own group ---
  \stickfigure{(-2.5, -1.9)}{SuperAdmin}
  \draw[link] (-1, -1.9) -- (uc1.west);
  \draw[link] (-1, -1.9) -- (uc2.west);
  \draw[link] (-1, -1.9) -- (uc3.west);
  \draw[link] (-1, -1.9) -- (uc4.west);

  \stickfigure{(-2.5, -10.0)}{Admin}
  \draw[link] (-1, -10.0) -- (uc5.west);
  \draw[link] (-1, -10.0) -- (uc6.west);
  \draw[link] (-1, -10.0) -- (uc7.west);
  \draw[link] (-1, -10.0) -- (uc8.west);
  \draw[link] (-1, -10.0) -- (uc9.west);
  \draw[link] (-1, -10.0) -- (uc10.west);
  \draw[link] (-1, -10.0) -- (uc11.west);
  \draw[link] (-1, -10.0) -- (uc12.west);
  \draw[link] (-1, -10.0) -- (uc13.west);

  \stickfigure{(-2.5, -16.25)}{Caissier}
  \draw[link] (-1, -16.25) -- (uc5.west);
  \draw[link] (-1, -16.25) -- (uc7.west);
  \draw[link] (-1, -16.25) -- (uc14.west);
  \draw[link] (-1, -16.25) -- (uc15.west);
  \draw[link] (-1, -16.25) -- (uc16.west);

  \stickfigure{(-2.5, -21.9)}{Membre}
  \draw[link] (-1, -21.9) -- (uc17.west);
  \draw[link] (-1, -21.9) -- (uc18.west);
  \draw[link] (-1, -21.9) -- (uc19.west);
  \draw[link] (-1, -21.9) -- (uc20.west);

  \end{tikzpicture}
  }%
  \caption{Diagramme de cas d'utilisation global de TrustGym}
  \label{fig:use_case_global}
\end{figure}

\subsection{Format des descriptions textuelles de cas d'utilisation}
\label{sec:uc_format}
Pour all\'{e}ger les chapitres de sprint, le format textuel des cas d'utilisation est d\'{e}fini une seule fois ici. Chaque sprint r\'{e}f\'{e}rence ce format et liste uniquement les variations sp\'{e}cifiques (acteur, sc\'{e}nario, exceptions notables) lorsqu'elles apportent une information non triviale.

\begin{itemize}
  \item \textbf{Acteur principal} : r\^{o}le qui d\'{e}clenche le cas d'utilisation (SuperAdmin, Admin, Caissier, Membre).
  \item \textbf{Pr\'{e}conditions} : utilisateur authentifi\'{e}, p\'{e}rim\`{e}tre tenant valide (partenaire/salle), permissions \acrshort{rbac} et \texttt{ModuleGuard} satisfaites.
  \item \textbf{Sc\'{e}nario nominal} : (1) l'acteur saisit les donn\'{e}es ou d\'{e}clenche l'action, (2) le client web/mobile valide localement et envoie la requ\^{e}te, (3) l'\acrshort{api} valide les \acrshort{dto} et les r\`{e}gles m\'{e}tier, (4) l'\acrshort{orm} persiste, (5) un \'{e}v\'{e}nement temps r\'{e}el est diffus\'{e} si pertinent, (6) le client met \`{a} jour son \'{e}tat.
  \item \textbf{Postcondition} : donn\'{e}es persist\'{e}es, journal d'audit ou \acrshort{ws} \'{e}mis, retour \texttt{2xx}.
  \item \textbf{Exceptions} : \texttt{401} non authentifi\'{e}, \texttt{403} permission/module manquant, \texttt{400} \acrshort{dto} invalide, \texttt{409} conflit m\'{e}tier (doublon, stock insuffisant, etc.).
\end{itemize}

\subsection{Product Backlog}
Le tableau~\ref{tab:product_backlog} liste les treize fonctionnalit\'{e}s du product backlog, tri\'{e}es par priorit\'{e} d\'{e}croissante (0--100). L'authentification et la gestion des membres arrivent en t\^{e}te --- aucune autre fonctionnalit\'{e} n'a de sens sans elles.

\begin{longtable}{|c|p{4.5cm}|p{6.5cm}|c|}
\hline
\textbf{ID} & \textbf{Titre} & \textbf{Description} & \textbf{Importance} \\
\hline
\endfirsthead
\hline
\textbf{ID} & \textbf{Titre} & \textbf{Description} & \textbf{Importance} \\
\hline
\endhead
\hline
\endfoot
1 & Authentification et gestion des utilisateurs & Mise en place du double syst\`{e}me \acrshort{jwt} (Admin/Membre), gestion des r\^{o}les et permissions, s\'{e}curisation des endpoints. & 100 \\
\hline
2 & Gestion des membres & \acrshort{crud} complet des membres, recherche, filtrage, consultation des profils et historiques. & 100 \\
\hline
3 & Gestion des abonnements & Cr\'{e}ation des types d'abonnements, affectation aux membres, suivi des \'{e}ch\'{e}ances et renouvellements. & 90 \\
\hline
4 & Contr\^{o}le d'acc\`{e}s & Gestion des entr\'{e}es/sorties, configuration des cr\'{e}neaux horaires, historique d'acc\`{e}s, calcul de l'occupation en temps r\'{e}el. & 90 \\
\hline
5 & Int\'{e}gration mat\'{e}rielle & Communication avec les terminaux ZKTeco et Hikvision via le module machine handler. & 90 \\
\hline
6 & Application mobile membre & Interface mobile compl\`{e}te : inscription multi-\'{e}tapes, tableau de bord, s\'{e}ances, statistiques, profil. & 90 \\
\hline
7 & Gestion des coachs & Ajout et gestion des coachs, affectation aux disciplines, gestion des disponibilit\'{e}s. & 80 \\
\hline
8 & Gestion des s\'{e}ances & Planification des s\'{e}ances, affectation des coachs, vue calendrier, gestion des inscriptions. & 80 \\
\hline
9 & Boutique et stock & Gestion des produits, des fournisseurs, suivi des stocks, enregistrement des ventes. & 70 \\
\hline
10 & Finance et facturation & G\'{e}n\'{e}ration de factures \acrshort{b2b}, suivi des paiements, gestion de la caisse, rapports financiers. & 70 \\
\hline
11 & Portefeuille pr\'{e}pay\'{e} & Syst\`{e}me de wallet pour les membres, rechargement, d\'{e}bit automatique, historique des transactions. & 60 \\
\hline
12 & Centre \acrshort{kpi} & Tableaux de bord analytiques, indicateurs de performance, graphiques et rapports. & 60 \\
\hline
13 & Assistant \acrshort{ai} (\acrshort{rag}) & Assistant intelligent utilisant OpenAI et pgvector pour r\'{e}pondre aux questions des administrateurs. & 50 \\
\hline
\caption{Product Backlog de TrustGym}
\label{tab:product_backlog}
\end{longtable}

\subsection{Diagramme de classes}
Le diagramme de classes (figure~\ref{fig:class_diagram}) capture la structure du mod\`{e}le de donn\'{e}es \cite{class_diag, fowler2003uml}. Il s'organise autour de quatre p\^{o}les : le noyau multi-tenant (\textbf{Partner}, \textbf{Gym}, \textbf{User}, \textbf{Member}), le contr\^{o}le d'acc\`{e}s (\textbf{Door}, \textbf{Access}, \textbf{GymSubscription}, \textbf{Payment}), l'op\'{e}rationnel m\'{e}tier (\textbf{Coach}, \textbf{Session}, \textbf{Product}, \textbf{Supplier}, \textbf{Invoice}, \textbf{Wallet}) et l'\acrshort{ai} (\textbf{AiChatLog}). Le catalogue exhaustif des 57 entit\'{e}s est donn\'{e} ci-dessous (tableau~\ref{tab:entity_catalog}).

\begin{figure}[H]
  \centering
  \includegraphics[width=\textwidth]{class diagram.png}
  \caption{Diagramme de classes global de TrustGym}
  \label{fig:class_diagram}
\end{figure}

\subsection{Catalogue complet des entit\'{e}s}

Le mod\`{e}le complet comprend \textbf{57 entit\'{e}s TypeORM} r\'{e}parties en neuf domaines. Le tableau~\ref{tab:entity_catalog} en donne le catalogue exhaustif.

\begin{longtable}{|c|>{\footnotesize}p{4.8cm}|p{7.5cm}|}
\hline
\textbf{\#} & \textbf{Entit\'{e}} & \textbf{Description} \\
\hline
\endfirsthead
\hline
\textbf{\#} & \textbf{Entit\'{e}} & \textbf{Description} \\
\hline
\endhead
\hline
\endfoot

\multicolumn{3}{|c|}{\textbf{Domaine : Noyau (Core)}} \\
\hline
1 & Partner & Partenaire \acrshort{saas} (franchise) ; informations fiscales, cotisation annuelle \\
\hline
2 & Gym & Salle de sport physique ; configuration bridge, cl\'{e} \acrshort{api}, heartbeat \\
\hline
3 & User & Utilisateur back-office (Admin, Caissier, SuperAdmin) ; cardId et PIN \\
\hline
4 & Role & R\^{o}les \acrshort{rbac} : SuperAdmin, Admin, Cashier, Coach, Member \\
\hline
5 & LoginLog & Journal d'audit des connexions/d\'{e}connexions \\
\hline
6 & RefreshToken & Suivi des tokens de rafra\^{i}chissement \\
\hline
7 & ResetPasswordOtp & Codes \acrshort{otp} pour la r\'{e}initialisation de mot de passe \\
\hline

\multicolumn{3}{|c|}{\textbf{Domaine : Membres}} \\
\hline
8 & Member & Adh\'{e}rent de la salle ; carte \acrshort{rfid}, PIN, mot de passe mobile \\
\hline
9 & MemberImage & R\'{e}f\'{e}rence S3 pour la photo du membre \\
\hline
10 & AccessMemberMember & Table de jointure pour les relations entre membres (famille) \\
\hline

\multicolumn{3}{|c|}{\textbf{Domaine : Abonnements et Paiements}} \\
\hline
11 & GymSubscription & Formule d'abonnement : prix, dur\'{e}e ou nombre de s\'{e}ances \\
\hline
12 & MemberGymSubscription & Instance d'abonnement : membre + formule + dates + paiements \\
\hline
13 & Payment & \'{E}ch\'{e}ance de paiement : montant, type, date, statut \\
\hline
14 & GymSubscriptionDoor & Liaison formule $\leftrightarrow$ porte autoris\'{e}e \\
\hline
15 & SubscriptionTimeSlot & Restrictions horaires pour les abonnements hors-pointe \\
\hline
16 & Convention & Groupe de r\'{e}duction (entreprise/famille) : montant fixe ou pourcentage \\
\hline
17 & MemberAnnualFee & Cotisation annuelle d'inscription par membre \\
\hline

\multicolumn{3}{|c|}{\textbf{Domaine : Contr\^{o}le d'Acc\`{e}s}} \\
\hline
18 & Door & Porte physique : IP, port, \acrshort{sdk}, direction (In/Out/InOut) \\
\hline
19 & Access & \'{E}v\'{e}nement d'entr\'{e}e/sortie : date, heure, autorisation, type \\
\hline
20 & DeviceCommand & File de commandes pour les dispositifs en mode PUSH \\
\hline
21 & DeviceSyncJob & Suivi des t\^{a}ches de synchronisation avec les dispositifs \\
\hline

\multicolumn{3}{|c|}{\textbf{Domaine : Coachs et S\'{e}ances}} \\
\hline
22 & Coach & Entra\^{i}neur : tarif horaire, carte d'acc\`{e}s, disponibilit\'{e}s \\
\hline
23 & CoachImage & R\'{e}f\'{e}rence S3 pour la photo du coach \\
\hline
24 & Discipline & Type d'activit\'{e} sportive (CrossFit, Boxe, Yoga, etc.) \\
\hline
25 & DisciplineImage & R\'{e}f\'{e}rence S3 pour l'image de la discipline \\
\hline
26 & Session & S\'{e}ance planifi\'{e}e : date, coach, discipline, type (OPEN/CLOSE/RESERVATION) \\
\hline
27 & SessionPayment & R\'{e}mun\'{e}ration du coach pour une s\'{e}ance \\
\hline
28 & SessionInvitation & R\'{e}servation d'un membre pour une s\'{e}ance \\
\hline
29 & SessionQRCode & Code \acrshort{qr} pour le pointage de pr\'{e}sence \\
\hline
30 & CoachTimeSlot & Cr\'{e}neaux de disponibilit\'{e} hebdomadaire du coach \\
\hline
31 & TimeSlotHours & Plages jour/heure sp\'{e}cifiques \\
\hline
32 & PartnerTimeSlot & Configuration horaire au niveau partenaire \\
\hline

\multicolumn{3}{|c|}{\textbf{Domaine : Boutique}} \\
\hline
33 & Product & Article en vente : nom, prix, cat\'{e}gorie \\
\hline
34 & ProductCategory & Classification des produits \\
\hline
35 & ProductImage & R\'{e}f\'{e}rence S3 pour l'image du produit \\
\hline
36 & ProductStock & Niveau d'inventaire par produit par salle \\
\hline
37 & MemberPurchase & Transaction de vente (POS) : membre achète un produit \\
\hline
38 & Supplier & Fournisseur de produits \\
\hline
39 & SupplierImage & R\'{e}f\'{e}rence S3 pour l'image du fournisseur \\
\hline
40 & SupplierPurchase & Commande d'achat aupr\`{e}s d'un fournisseur \\
\hline
41 & StockTransfer & Mouvement d'inventaire inter-salles \\
\hline
42 & StockMovement & Piste d'audit de tous les changements de stock \\
\hline

\multicolumn{3}{|c|}{\textbf{Domaine : Finance}} \\
\hline
43 & Wallet & Solde pr\'{e}pay\'{e} par membre \\
\hline
44 & WalletEntry & Transaction individuelle du portefeuille (cr\'{e}dit/d\'{e}bit) \\
\hline
45 & Invoice & Facture \acrshort{b2b} avec champs fiscaux tunisiens \\
\hline
46 & InvoiceLine & Ligne de d\'{e}tail sur une facture \\
\hline
47 & CashRegister & Solde de caisse par salle \\
\hline
48 & CashRegisterLog & Journal des transactions de caisse \\
\hline

\multicolumn{3}{|c|}{\textbf{Domaine : \acrshort{saas} et Plateforme}} \\
\hline
49 & GymPartnerSubscription & Abonnement \acrshort{saas} de la plateforme par partenaire \\
\hline
50 & PartnerSubscriptionType & D\'{e}finition des paliers (Core, Professional, Business, Enterprise) \\
\hline
51 & PartnerSubscriptionModule & Modules de fonctionnalit\'{e}s par palier d'abonnement \\
\hline

\multicolumn{3}{|c|}{\textbf{Domaine : Intelligence Artificielle}} \\
\hline
52 & AiChatLog & Historique des conversations \acrshort{ai} \\
\hline
53 & DocChunk & Morceaux de documents avec embeddings vectoriels pour le \acrshort{rag} \\
\hline
54 & AiActions & Suivi des actions de l'assistant \acrshort{ai} \\
\hline

\multicolumn{3}{|c|}{\textbf{Tables de jointure}} \\
\hline
55--57 & Tables de jointure & Relations ManyToMany (Gym$\leftrightarrow$Subscription, etc.) \\
\hline

\caption{Catalogue des 57 entit\'{e}s de la base de donn\'{e}es TrustGym}
\label{tab:entity_catalog}
\end{longtable}

\section{Conclusion}
Ce chapitre a \'{e}tabli le contrat fonctionnel du projet : qui utilise TrustGym, ce que chacun peut faire, et les contraintes qualit\'{e} attendues. Le product backlog a fix\'{e} les priorit\'{e}s, le mod\`{e}le de donn\'{e}es a pos\'{e} les fondations. Les chapitres suivants d\'{e}crivent comment tout cela a \'{e}t\'{e} construit, sprint par sprint.

\end{document}
